\section{Morphisms of arrow-based conditions}
\slabel{ab-morphisms}

The notion of \emph{morphism} of nested conditions has not received much attention in the literature. When considered at all, such as in \cite{bchk:conditional-reactive-systems,sksclo:coinductive-techniques-for-satisfiability}, morphisms are essentially based on the semantics in terms of satisfaction. Indeed, entailment (see \cref{def:ab-satisfaction}) establishes a preorder over conditions. Let us denote the resulting category $\ABC^{\entails}$, where if $c$ and $b$ are ab-conditions, $c \leq b$ iff $c \entails b$.%
%: Note that in this case $c$ and $b$ have the same root. %
\footnote{Note that $c \leq b$ implies that $c$ and $b$ have the same root.} 
 
The question that we address in this section is to establish a meaningful \emph{structural} notion of condition morphism. That is, given the fact that an ab-condition is essentially a diagram in the category $\bC$, a structural morphism from $b$ to $c$ consists of arrows between objects in $b$ and $c$ (when viewed as diagrams), satisfying certain conditions on the resulting subdiagrams. 

For morphisms to be meaningful,  they should only exist where there is entailment: 
% indeed, we require there to be a \emph{contravariant} functor to $\ABC^{\entails}$, meaning that 
if there is a structural morphism $b \to c$, then it should hold $c \entails b$.\footnote{This is consistent with the case of the $\exists$-fragment sketched in the Introduction.} As we will see, this requirement implies that it is not enough to consider only arrows \emph{from} objects of $b$ \emph{to} objects of $c$, since in our definition of satisfaction the direction of entailment flips at every subsequent nesting level. 


\todo[inline]{From here is the version based on root arrows.}

The underlying principle for a morphism $m$ from $b$ to $c$ is that it must identify, for every $c$-branch $p^c_i$, a $b$-branch $p^b_j$ that it entails. However, because of their inductively defined tree-shape, on lower levels we have to relate (sub)conditions over different roots; hence, in general $R^c$ and $R^b$ are not the same. Instead, a morphism $m$ will include an arrow $v\of R^c\to R^b$ (in the \emph{opposite} direction of $m$), and the entailment of $p^b_j$ by $p^c_i$ is \emph{modulo} $v$, in a way made precise below.\todo{AR: Reflection square? Reflection diagram? ab-reflection diagram?}

\begin{definition}[reflection square]\label{def:reflection square}
Let $a_1\of R_1\to P_1$, $a_2\of R_2\to P_2$, $v\of R_2\to R_1$ and $v'\of P_1\to P_2$, and consider the (not necessarily commuting) diagram on the left:
\begin{equation}\label{eq:reflection square}
\begin{tikzpicture}[on grid,node distance = 1.5cm,baseline=(a1)]
\node (R1) {$R_1$};
\node (P1) [below=of R1] {$P_1$};
\node (R2) [right=of R1] {$R_2$};
\node (P2) [below=of R2] {$P_2$};

\path
  (R1) edge [->] node (a1) [auto] {$a_1$} (P1)
  (R2) edge [->] node [auto] {$a_2$} (P2)
       edge [->] node [above] {$v$} (R1)
  (P1) edge [->] node [auto] {$v'$} (P2);
\end{tikzpicture}
\qquad
\begin{tikzpicture}[on grid,node distance = 1.5cm,baseline=(a1)]
\node (R1) {$R_1$};
\node (P1) [below=of R1] {$P_1$};
\node (R2) [right=of R1] {$R_2$};
\node (P2) [below=of R2] {$P_2$};

\path
  (R1) edge [->] node (a1) [auto] {$a_1$} (P1)
  (R2) edge [->] node [auto] {$a_2$} (P2)
       edge [->] node [above] {$v$} (R1)
  (P1) edge [->] node [auto] {$v'$} (P2);

\node (G) [below left=of P1] {$G$};
\path
  (R1) edge [->,bend right=25] node[left] {$g$} (G)
  (P2) edge [->,bend left=25] node[auto] {$h$} (G);
  
\path (G) edge[pushout] +(+3.5mm,3.5mm);
\path (G) edge[white] node[pos=.7,black] {$\circlearrowright$} (P1);
\end{tikzpicture}
\qquad
\begin{tikzpicture}[on grid,auto,node distance = 1.5cm,baseline=(a1)]
\node (R1) {$R_1$};
\node (P1) [below=of R1] {$P_1$};
\node (R2) [right=of R1] {$R_2$};
\node (P2) [below=of R2] {$P_2$};

\path
  (R1) edge [->] node (a1) [auto] {$a_1$} (P1)
  (R2) edge [->] node [auto] {$a_2$} (P2)
       edge [->] node [above] {$v$} (R1)
  (P1) edge [->] node [auto] {$v'$} (P2);

\node (E) [right=of R2] {$E$};
\path (E) edge [->] node[above] {$e$} (R2);
\end{tikzpicture}
\end{equation}
%
\begin{itemize}[defsep]
\item The diagram is called an \emph{arrow-based reflection square} if, for the pushout $g,h$ of the span $v,a_2$ as drawn in the middle diagram, the bottom left subdiagram commutes, i.e., $g=a_1;v';h$.
\item The diagram is called a \emph{simple arrow-based reflection square} if there is an equalising arrow $e$ as in the right hand diagram (i.e., such that $e;a_2=e;v;a_1v'$) such that $e;v$ is epi.
\end{itemize}
\end{definition}
%
The following proposition states some useful facts regarding [simple] reflection squares.

\begin{proposition}[reflection squares]\label{prop:reflection square}~
\begin{enumerate}[defsep]
\item Every simple reflection square is a reflection square.

\item In the left hand diagram of \eqref{eq:reflection square}, if $v$ is epi and the diagram commutes (i.e., $a_2=v;a_1;v'$), it is a simple reflection square.

\item\label{simple reflection square equaliser} In a simple reflection square (right hand diagram of \eqref{eq:reflection square}), if $e$ is the equaliser of $a_2$ and $v;a_1;v'$ then $e;v$ is epi.

\item In the left hand diagram of \eqref{eq:reflection square}, if $v$ is split epi with section $s$ such that $s;a_2=a_1;v'$, the diagram is a simple reflection square.

\begin{equation}\label{eq:reflection square section}
\begin{tikzpicture}[on grid,node distance = 1.5cm,baseline=(a1)]
\node (R1) {$R_1$};
\node (P1) [below=2 of R1] {$P_1$};
\node (R2) [right=3 of R1] {$R_2$};
\node (P2) [below=2 of R2] {$P_2$};

\path
  (R1) edge [->] node [left] (a1) {$a_1$} (P1)
  (R2) edge [->] node [auto] {$a_2$} (P2)
       edge [->] node [above] (v) {$v$} (R1)
  (P1) edge [->] node [below] (v') {$v'$} (P2);
  
\node (S) [below right=1 and 1.5 of R1] {$R_1$};

\path (S) edge[->] node[below left,pos=.4,inner sep=3] {$\id$} (R1)
          edge[->] node[below right,pos=.4,inner sep=3] {$s$} (R2);
          
\path (S) edge[white] node[black] {$\circlearrowright$} (v')
          edge[white] node[black] {$\circlearrowright$} (v);
\end{tikzpicture}
\end{equation}

\item Simple reflection squares compose; i.e., in
%
\begin{equation}\label{eq:reflection squares compose}
\begin{tikzpicture}[on grid,node distance = 1.5cm,baseline=(a1)]
\node (R1) {$R_1$};
\node (P1) [below=of R1] {$P_1$};
\node (R2) [right=of R1] {$R_2$};
\node (P2) [below=of R2] {$P_2$};
\node (R3) [right=of R2] {$R_3$};
\node (P3) [below=of R3] {$P_3$};

\path
  (R1) edge [->] node (a1) [auto] {$a_1$} (P1)
  (R2) edge [->] node [auto] {$a_2$} (P2)
       edge [->] node [above] {$v_1$} (R1)
  (P1) edge [->] node [auto] {$v_1'$} (P2)
  (R3) edge [->] node [auto] {$a_3$} (P3)
       edge [->] node [above] {$v_2$} (R2)
  (P2) edge [->] node [auto] {$v_2'$} (P3);
\end{tikzpicture}
\end{equation}
%
if the left and right hand subdiagrams are simple reflection squares, then so is the outer subdiagram.

\item In a reflection square, if the upper arrow $v$ is an isomorphism, then the diagram commutes.

\item\label{reflection square any} In a reflection square, for \emph{any} pair $g',h'$ that makes the outer subdiagram commute (i.e., such that $v;g'=a_2;h'$), the bottom left subdiagram commutes (i.e., $g'=a_1;v';h'$):
\begin{equation}\label{eq:reflection square any}
\begin{tikzpicture}[on grid,node distance = 1.5cm,baseline=(a1)]
\node (R1) {$R_1$};
\node (P1) [below=of R1] {$P_1$};
\node (R2) [right=of R1] {$R_2$};
\node (P2) [below=of R2] {$P_2$};

\path
  (R1) edge [->] node (a1) [auto] {$a_1$} (P1)
  (R2) edge [->] node [auto] {$a_2$} (P2)
       edge [->] node [above] {$v$} (R1)
  (P1) edge [->] node [auto] {$v'$} (P2);

\node (G) [below left=of P1] {$G'$};
\path
  (R1) edge [->,bend right=25] node[left] {$g'$} (G)
  (P2) edge [->,bend left=25] node[auto] {$h'$} (G);
  
\path (G) edge[white] node[pos=.7,black] {$\circlearrowright$} (P1);
\end{tikzpicture}
\end{equation}

\item Reflection squares compose; i.e., in \eqref{eq:reflection squares compose}, 
if both the left hand and the right hand subdiagrams are reflection squares, then so is the outer subdiagram.
\end{enumerate}
\end{proposition}

\begin{proof}[]
\begin{enumerate}[defsep]
\item The equalising arrow guarantees $e;a_2;h=e;v;a_1;v';h$, and the pushout square guarantees $e;a_2;h=s;v;g$; hence $e;v;a_1;v';h=e;v;g$. Since $e;v$ is epi, it follows that $a_1;v';h=g$.

\item Since the diagram already commutes, $e=\id$ is the equaliser; hence $e;v=v$ is epi.

\item In the right hand diagram of \eqref{eq:reflection square}, let $e$ be such that $e;a_2=e;v;a_1;v'$ and $e;v$ is epi, and consider the equaliser $e'\of E'\to R_2$ of $a_2$ and $v;a_1;v'$. By the universal property of $e'$, there is a (unique) $k\of E\to E'$ such that $e=k;e'$. But then, since $e;v$ ($=k;e';v$) is epi, so is $e';v$.

\[\begin{tikzpicture}[on grid,auto,node distance = 1.5cm,baseline=(a1)]
\node (R1) {$R_1$};
\node (P1) [below=of R1] {$P_1$};
\node (R2) [right=of R1] {$R_2$};
\node (P2) [below=of R2] {$P_2$};

\path
  (R1) edge [->] node (a1) [left] {$a_1$} (P1)
  (R2) edge [->] node [left] {$a_2$} (P2)
       edge [->] node [above] {$v$} (R1)
  (P1) edge [->] node [auto] {$v'$} (P2);

\node (E) [right=of R2] {$E$};
\path (E) edge [->] node[above] {$e$} (R2);

\node (E') [below=of E] {$E'$};
\path (E') edge [right hook->] node[pos=.4,inner sep=1] {$e'$} (R2)
      (E) edge [dashed,->] node[right] {$k$} (E');
\end{tikzpicture}\]

\item The section $s$ in \eqref{eq:reflection square section} satisfies the conditions for $e$ in the right hand diagram of \eqref{eq:reflection square} (noting that $s;v=\id$ is epi).

\item Consider the following diagram, where $e_1$ is such that \tagone $e_1;a_2=e_1;v_1;a_1;v_1'$ and $e_1;v_1$ is epi, $e_2$ is such that \tagtwo $e_2;a_3=e_2;v_2;a_2;v_2'$ and $e_2;v_2$ is epi, and $e_1',e_2'$ form the pullback of $e_1$ and $e_2;v_2$.
%
\[
\begin{tikzpicture}[on grid,auto,node distance = 1.5cm,baseline=(a1)]
\node (R1) {$R_1$};
\node (P1) [below=of R1] {$P_1$};
\node (R2) [right=of R1] {$R_2$};
\node (P2) [below=of R2] {$P_2$};
\node (R3) [right=of R2] {$R_3$};
\node (P3) [below=of R3] {$P_3$};

\path
  (R1) edge [->] node (a1) [auto] {$a_1$} (P1)
  (R2) edge [->] node [auto] {$a_2$} (P2)
       edge [->] node [above] {$v_1$} (R1)
  (P1) edge [->] node [auto] {$v_1'$} (P2)
  (R3) edge [->] node [auto] {$a_3$} (P3)
       edge [->] node [above] {$v_2$} (R2)
  (P2) edge [->] node [auto] {$v_2'$} (P3);
  
\node (E1) [above=of R2] {$E_1$};
\node (E2) [right=of R3] {$E_2$};
\node (E) [above=of E2] {$E$};

\path
  (E) edge [->] node {$e_1'$} (E2)
      edge [->] node[above] {$e_2'$} (E1)
  (E1) edge [left hook->] node {$e_1$} (R2)
  (E2) edge [left hook->] node[above] {$e_2$} (R3);

\path (E) edge[pushout] +(-3.5mm,-3.5mm);
\end{tikzpicture}
\]
%
Using also that \tagthree the pullback square commutes, we can deduce
\begin{align*}
e'_1;e_2;v_2;v_1;a_1;v_1';v_2'
 & \stackrel{\tagthree}{=} e_2';e_1;v_1;a_1;v_1';v_2' \\
 & \stackrel{\tagone}= e_2';e_1;a_2;v_2' \\
 & \stackrel{\tagthree}= e_1';e_2;v_2;a_2;v_2' \\
 & \stackrel{\tagtwo}= e_1';e_2;a_3
\end{align*}
%
Moreover, epis are preserved by pullback (in presheaf toposes), hence $e_2'$ is epi. Finally, epis compose, hence $e_1';e_2;v_2;v_1=e_2';e_1;v_1$ is epi. Thus, $e_1';e_2$ satisfies the conditions for $e$ in the definition of simple reflection squares..

\item Consider the middle diagram of \eqref{eq:reflection square} in case $v$ is an isomorphism. It follows that $h$ is also iso; hence $v;g=a_2;h$ and $g=a_1;v';h$ can be equivalently rewritten to \tagone $v;g;h^{-1}=a_2$ and \tagtwo $g;h^{-1}=a_1; v'$, respectively. Substituting \tagtwo into \tagone, we obtain $v;a_1;v'=a_2$, hence the reflection square commutes.


\item For a pair $g', h'$ as in \eqref{eq:reflection square any}, there exists a (unique) $k\of G\to G'$ (where $G$ is the pushout object as in \eqref{eq:reflection square}) such that $g'=g;k$ and $h'=h;k$; hence $g'=a_1; v';h;k=a_1;v';h'$.

\item Consider the pushout $g,h$ of $v_1;v_2$ and $a_3$:

\[\begin{tikzpicture}[auto,on grid,node distance = 1.5cm,baseline=(a1)]
\node (R1) {$R_1$};
\node (P1) [below=of R1] {$P_1$};
\node (R2) [right=of R1] {$R_2$};
\node (P2) [below=of R2] {$P_2$};
\node (R3) [right=of R2] {$R_3$};
\node (P3) [below=of R3] {$P_3$};

\path
  (R1) edge [->] node (a1) [auto] {$a_1$} (P1)
  (R2) edge [->] node [auto] {$a_2$} (P2)
       edge [->] node [above] {$v_1$} (R1)
  (P1) edge [->] node [auto] {$v_1'$} (P2)
  (R3) edge [->] node [auto] {$a_3$} (P3)
       edge [->] node [above] {$v_2$} (R2)
  (P2) edge [->] node [auto] {$v_2'$} (P3);

\node (G) [below left=of P1] {$G$};
\path
  (R1) edge [->,bend right=25] node[left] {$g$} (G)
  (P3) edge [->,bend left=20] node[auto] {$h$} (G);
  
\path (G) edge[pushout] +(.4,.4);
\end{tikzpicture}\]
%
Due to Clause~\ref{reflection square any} applied to the right hand subdiagram, $v_1;g=a_2;v_2';h$; but then due to the same clause applied to the left hand subdiagram, $g=a_1;v'_1;v_2';h$.
\qed
\end{enumerate}
\end{proof}
%
It is an immediate observation that, for all reflection squares obtained by applying Prop.~\ref{prop:reflection square}, the upper arrow $v$ is epi; however, we currently have no proof whether that is a necessary condition.\todo{mathexchange}

The dual to a reflection square is a \emph{preservation square}.

\begin{definition}[preservation square]\label{def:preservation square}
Let $a_1\of R_1\to P_1$, $a_2\of R_2\to P_2$, $v\of R_2\to R_1$ and $v'\of P_1\to P_2$, and consider the (not necessarily commuting) diagram on the left:
\begin{equation}\label{eq:preservation square}
\begin{tikzpicture}[on grid,node distance = 1.5cm,baseline=(a1)]
\node (R1) {$R_1$};
\node (P1) [below=2 of R1] {$P_1$};
\node (R2) [right=2 of R1] {$R_2$};
\node (P2) [below=2 of R2] {$P_2$};

\path
  (R1) edge [->] node (a1) [auto] {$a_1$} (P1)
  (R2) edge [->] node [auto] {$a_2$} (P2)
       edge [->] node [above] {$v$} (R1)
  (P1) edge [->] node [auto] {$v'$} (P2);
\end{tikzpicture}
\qquad\qquad
\begin{tikzpicture}[on grid,node distance = 1.5cm,baseline=(a1)]
\node (R1) {$R_1$};
\node (P1) [below=2 of R1] {$P_1$};
\node (R2) [right=3 of R1] {$R_2$};
\node (P2) [below=2 of R2] {$P_2$};

\path
  (R1) edge [->] node [left] (a1) {$a_1$} (P1)
  (R2) edge [->] node [auto] {$a_2$} (P2)
       edge [->] node [above] (v) {$v$} (R1)
  (P1) edge [->] node [below] (v') {$v'$} (P2);
  
\node (S) [above right=1 and 1.5 of P1] {$P_1$};

\path (S) edge[<-] node[above left,pos=.4,inner sep=3] {$\id$} (P1)
          edge[<-] node[above right,pos=.4,inner sep=3] {$r'$} (P2);
          
\path (S) edge[white] node[black,pos=.4] {$\circlearrowright$} (v')
          edge[white] node[black] {$\circlearrowright$} (v);
\end{tikzpicture}
\end{equation}
%
The diagram is called an \emph{arrow-based preservation square} if $v'$ is a split mono with a retract $r'$ (meaning $v';r'=\id$) such that the upper pentagon in the right hand diagram above commutes (meaning $a_2;r'=v;a_1$).
\end{definition}
%
The following proposition states some useful facts regarding preservation squares.
%
\begin{proposition}[preservation squares]\label{prop:preservation square}~
\begin{enumerate}[defsep]
\item In the left hand diagram of \eqref{eq:preservation square}, if the bottom arrow $v'$ is iso then the diagram is a preservation square if and only if it commutes.

\item\label{preservation square any} In the left hand diagram of \eqref{eq:preservation square}, if $v$ is iso then the diagram is a preservation square if and only if, for any pair of arrows $g,h$ such that $g=a_1;h$, there is an arrow $h'$ such that $h=v';h'$ and $v;g=a_2;h'$.

\begin{equation}\label{eq:preservation square any}
\begin{tikzpicture}[on grid,node distance = 1.5cm,baseline=(a1)]
\node (R1) {$R_1$};
\node (P1) [below=of R1] {$P_1$};
\node (R2) [right=of R1] {$R_2$};
\node (P2) [below=of R2] {$P_2$};

\path
  (R1) edge [->] node (a1) [auto] {$a_1$} (P1)
  (R2) edge [->] node [auto] {$a_2$} (P2)
       edge [->] node [above] {$v$} (R1)
  (P1) edge [->] node [auto] {$v'$} (P2);

\node (G) [below left=of P1] {$G$};
\path (R1) edge[->,bend right=25] node[auto] (g) [left] {$g$} (G)
      (P1) edge[->] node[below right,inner sep=2,pos=.4] {$h$} (G)
      (P2) edge[dashed,->,bend left=25] node[auto] (h') {$h'$} (G);
      
\path (P1) edge[white] node[black,pos=.4] {$\circlearrowright$} (g)
           edge[white] node[black,pos=.4] {$\circlearrowright$} (h');
\end{tikzpicture}
\end{equation}

\item\label{preservation square compose} Preservation squares compose; i.e., in
%
\begin{equation}\label{eq:preservation squares compose}
\begin{tikzpicture}[on grid,node distance = 1.5cm,baseline=(a1)]
\node (R1) {$R_1$};
\node (P1) [below=of R1] {$P_1$};
\node (R2) [right=of R1] {$R_2$};
\node (P2) [below=of R2] {$P_2$};
\node (R3) [right=of R2] {$R_3$};
\node (P3) [below=of R3] {$P_3$};

\path
  (R1) edge [->] node (a1) [auto] {$a_1$} (P1)
  (R2) edge [->] node [auto] {$a_2$} (P2)
       edge [->] node [above] {$v_1$} (R1)
  (P1) edge [->] node [auto] {$v_1'$} (P2)
  (R3) edge [->] node [auto] {$a_3$} (P3)
       edge [->] node [above] {$v_2$} (R2)
  (P2) edge [->] node [auto] {$v_2'$} (P3);
\end{tikzpicture}
\end{equation}
%
if the left and right hand subdiagrams are preservation squares, then so is the outer subdiagram.
\end{enumerate}
\end{proposition}

\begin{proof}[]
\begin{enumerate}[defsep]
\item If $v'$ is iso in the left hand diagram of \eqref{eq:preservation square}, then it is split mono with retract $r'=v'^{-1}$. If the diagram commutes, then $a_2=v;a_1;v'$ and hence $a_2;r'=v_1;a_1$, proving the preservation square property. Vice versa, $a_2;r'=v_1;a_1$ implies $a_2=a_2;r';v'=v_1;a_2;v'$, proving commutativity of the original square.

\item For the \emph{if} part, note that $g=a_1$ and $h=\id$ are a particular pair of arrows such that $g=a_1;h$; hence the $h'$ such that $h'=v';h$ and $v;g=a_2;h$ is precisely the retract $r'$ required in the definition of preservation squares. For the \emph{only if} part, assume that $r'$ is the retract guaranteed by the preservation square and define $h'=r';h$; then $v';h'=v';r';h=h$ and $v;g=v;a_1;h=a_2;r';h=a_2;h'$ as required.

\item If both subdiagrams are preservation squares and the split monos $v_i'$ have retracts $r_i'$ for $i=1,2$, then $v_1';v_2'$ is also split mono, with retract $r_2';r_1'$. By applying $v_i;a_i=a_{i+1};r_i'$ for $i=1,2$, it follows that $v_2;v_1;a_1=v_2;a_2;r_1'=a_3;r_2';r_1'$.\qed
\end{enumerate}
\end{proof}

\begin{definition}[arrow-based condition morphism]\label{def:ab-morphism}
\begin{itemize}[defsep]
\item Let $c_i = (R_i,p^i_1\ccdots p^i_{w_i})\in \AC{R_i}$ for $i=1,2$. An \emph{ab-condition reflection [preservation] morphism} (or just ab-reflection [-preservation] morphism) $m\of c_1 \to c_2$ is a tuple $(v,o,\setof{m_{j,k}}_{(j,k)\in o})$ such that
\begin{itemize}[nosep]
\item $v\of R_2\to R_1$ is an arrow;
\item $o\subseteq [1,w_1]\times [1,w_2]$ is relation between $c_1$-branches and $c_2$-branches that is surjective and injective [total and univalent];
%\item $o\of [1,w_2]\to [1,w_1]$ [$o\of [1,w_1]\to [1,w_2]$] is a mapping from $c_2$-branches to $c_1$-branches [from $c_1$-branches to $c_2$-branches];
\item For all $(j,k)\in o$, $(v,m_{j,k})\of p^1_j\to p^2_k$ is an ab-branch reflection [preservation] morphism.
\end{itemize}
%
The component $v$ of an ab-condition morphism $m$ is called its \emph{root arrow} and is denoted $v_m$. $m$ is called \emph{root-preserving} if $R_1=R_2=R$ and $v_m=\id_R$.

\item Let $p_i=(a_i,c_i)\in \AB{R_i}$ for $i=1,2$. An \emph{ab-branch reflection [preservation] morphism} from $p_1$ to $p_2$ is a tuple $(v,m)\of p_1\to p_2$ such that
\begin{itemize}[nosep]
\item $v\of R_2\to R_1$ is an arrow;
\item $m\of c_2\to c_1$ is an ab-condition reflection [preservation] morphism;
\item $a_1$ and $a_2$ together with $v$ and $v_m$ form a reflection [preservation] square.
\end{itemize}
\end{itemize}
\end{definition}
%
Note that if $\depth(c_1)=0$ or $\depth(c_2)=0$ for the source $c_1$ or the target $c_2$ of some ab-condition morphism $m\of c_1\to c_2$, $o$ is empty, hence the requirement on the branch morphisms is vacuously fulfilled.

For morphism composition, it is important to realise that if two binary relations $\rho_1\subseteq X_1\times X_2$ and $\rho_2\subseteq X_2\times X_3$ are either both injective and surjective (as the $o$ in an ab-reflection morphism), or both total and univalent (as the $o$ in an ab-preservation morphism), then for any $(x_1,x_3)\in \rho_1;\rho_2$ there is exactly one $x_2\in X_2$ such that $(x_1,x_2)\in \rho_1$ and $(x_2,x_3)\in \rho_2$.

\begin{definition}[identities and composition of ab-morphisms]\label{def:ab-reflection composition}
Let $c_i=(R_i,p^i_1\ccdots p^i_{w_i})\in \AC{R_i}$ be ab-conditions for $i=1,2,3$ and let $m_i=(v_i,o_i,\setof{m^i_{j,k}}_{(j,k)\in o_i}) \of c_i\to c_{i+1}$ be ab-reflection morphisms for $i=1,2$. Then $\id_{c_1}$ and $m_1;m_2$ are defined as follows (where $c^1_j$ is the subcondition of branch $p^1_j$ for $j\in [1,w_1]$):
%
\begin{eqnarray}
\id_{c_1} & =
  & (\id_{R_1},\id_{[1,w_1]},\setof{\id_{c^1_(j,j)}}_{1\leq j\leq w_1})
  \label{eq:id condition morphism} \\
m_1;m_2 & =
  & (v_2;v_1,\, o_1;o_2,\, \setof{m^2_{l,k};m^1_{j,l}}_{(j,l)\in o_1,(l,k)\in o_2)}) \enspace.
 \label{eq:condition morphism composition}
\end{eqnarray}
\end{definition}

\begin{proposition}[ab-reflection morphisms reflect satisfaction]
  \label{preop:ab-reflection morphisms}
Let $b\in \AC R,c \in \AC{S}$ be arrow-based conditions. If $m:b\to c$ is an ab-reflection morphism, then $v_m;g\sat c$ implies $g \sat b$ for all arrows $g\of R\to G$. 
\end{proposition}

\begin{proof}[]
By induction on $\depth(b)+\depth(c)$. Let $b=(R,p_1\cdots p_{|b|})$ such that $p_j=(a^p_j,c^p_j)$ for all $j\in[1,|b|]$ and $c=(S,q_1\cdots q_{|c|})$ such that $q_k=(a^q_k,c^q_k)$ for all $k\in[1,|c|]$, and let $m=(v,o,\setof{m_{j,k}}_{(j,k)\in o})$.
\begin{description}
\item[Base case] If $\depth(b)+\depth(c)=0$, it follows that $|c|=0$; hence $v_m;g\sat c$ is contradictory.
\item[Induction step.] Assume $v;g\sat c$, due to $v;g\sat q_k$ for some $k\in [1,|c|]$; and let $h$ be the witness, meaning $v;g=a^q_k;h$ and $h\not\sat c^q_k$. Let $(j,k)\in o$. Because $(v,m_{j,k})\of p_j\to q_k$ is an ab-branch morphism, $a^p_j$ and $a^q_k$ together with $v$ and $v_{m_{j,k}}$ form a reflection square, hence $g=a^p_j;v_{m_i};h$.

\smallskip
Recall that $m_{j,k}\of c^q_k\to c^p_j$ is an ab-reflection morphism and observe that $\depth(c^q_k)+\depth(c^p_j) < \depth(b)+\depth(c)$; hence by the induction hypothesis, $m_{j,k}$ fulfils the property. Therefore, if $v_{m_i};h\sat c^p_j$ then $h\sat c^q_k$, which is contradictory; it follows that $v_{m_{j,k}};h\not\sat c^p_j$, from which we may conclude $g\sat p_j$ and hence $g\sat b$. \qed
\end{description}
\end{proof}

\begin{proposition}[ab-preservation morphisms preserve satisfaction]
  \label{prop:ab-preservation morphisms preserve}
Let $b\in \AC R,c \in \AC{S}$ be arrow-based conditions. If $m:b\to c$ is an ab-preservation morphism, then $g\sat b$ implies $v_m;g \sat c$ for all arrows $g\of R\to G$. 
\end{proposition}

\begin{proof}[]
By induction on $\depth(b)+\depth(c)$. Let $b=(R,p_1\cdots p_{|b|})$ such that $p_j=(a^p_j,c^p_j)$ for all $j\in[1,|b|]$ and $c=(S,q_1\cdots q_{|c|})$ such that $q_k=(a^q_k,c^q_k)$ for all $i\in[1,|c|]$, and let $m=(v,o,\setof{m_{j,k}}_{(j,k)\in o})$.
\begin{description}
\item[Base case] If $\depth(b)+\depth(c)=0$, it follows that $|b|=0$; hence $g\sat b$ is contradictory.
\item[Induction step.] Assume $g\sat b$, due to $g\sat p_j$ for some $j\in [1,|b|]$; and let $h$ be the witness, meaning $g=a^p_j;h$ and $h\not\sat c^p_j$. Let $(j,k)\in o$. Because $(v,m_{j,k})\of p_{i}\to q_{o(i)}$ is an ab-branch preservation morphism, $a^p_j$ and $a^q_k$ together with $v$ and $v_{m_{j,k}}$ form a preservation square, hence (due to \cref{prop:preservation square}.\ref{preservation square any}) there is a $h'$ such that $h=v_{m_{j,k}};h'$ and $v_m;g=a^q_k;h'$.

\smallskip
Recall that $m_i\of c^p_k\to c^q_j$ is an ab-preservation morphism and observe that $\depth(c^q_k)+\depth(c^p_j) < \depth(b)+\depth(c)$; hence by the induction hypothesis, $m_{j,k}$ fulfils the property. Therefore, if $h'\sat c^q_k$ then $h=v_{m_{j,k}};h'\sat c^p_j$, which is contradictory; it follows that $h'\not\sat c^q_k$, from which we may conclude $v_m;g\sat q_k$ and hence $v_m;g\sat c$. \qed
\end{description}
\end{proof}
%
Thus, both reflection morphisms and preservation morphisms provide proof techniques for entailment: in particular, every root-preserving reflection morphism $m\of b\to c$ as well as  every root-preserving preservation morphism $m\of c\to b$ constitutes evidence for $c\entails b$. However, it turns out that reflection morphisms are strictly more powerful than preservation morphisms, in the following sense.
%
\begin{proposition}\label{prop:reflection beats preservation}
If $m\of b\to c$ is a preservation morphism of which $v_m$ is split mono, then there is a reflection morphism $m'\of c\to b$ such that, moreover, $m;m'=\id_b$.
\end{proposition}
%
\begin{proof}[]
By induction on $\depth(b)+\depth(c)$.
\begin{description}[defsep]
\item[Base case.] If $\depth(b)+\depth(c)=0$, then $m=(v,\varnothing,\varnothing)$ where $v$ is split mono; hence there is a retract $r\of R^c\to R^b$ such that $v;r=\id_{R^b}$. Let $m'=(r,\varnothing,\varnothing)$. Clearly $m'$ is a morphism from $c$ to $b$ and $m;m'=\id_b$.

\item[Induction step.] Assume the property holds up to $\depth(b)+\depth(c)$, and let $b=(R,p_1\ccdots p_{w_b})$, $c=(S,q_1\ccdots q_{w_c})$ and  $m=(v,o, \setof{m_{j,k}}_{(j,k)\in o})$ such that $o$ is injective and surjective.

\smallskip
For all $j\in[1,w_b]$ and $k\in [1,w_c]$, because $v$ is a split mono with retract $r$ and $(v,m'_{j,k})\of p_j\to q_k$ is a branch preservation morphism, the diagram on the left exists, in which all subdiagrams commute:
%
\[\begin{tikzpicture}[on grid,node distance = 1.5cm,baseline=(a1)]
\node (R1) {$R^b$};
\node (P1) [below=2 of R1] {$P^p_j$};
\node (R2) [right=3 of R1] {$R^c$};
\node (P2) [below=2 of R2] {$P^q_k$};

\path
  (R1) edge [->] node (a1) [left] {$a^p_j$} (P1)
  (R2) edge [->] node [auto] {$a^q_k$} (P2)
       edge [->] node [below] {$v$} (R1)
  (P1) edge [->] node [below] (v') {$v_{m_{j,k}}$} (P2);

\node (S') [above right=1 and 1.5 of P1] {$P^p_j$};  

\path (S') edge[<-] node[above left,pos=.4,inner sep=2] {$\id$} (P1)
          edge[<-] node[above right,pos=.4,inner sep=2] {$r_{k,j}$} (P2);

\node (S) [above right=1 and 1.5 of R1] {$R^c$};  

\path (S) edge[<-] node[above left,pos=.4,inner sep=3] {$r$} (R1)
          edge[<-] node[above right,pos=.4,inner sep=3] {$\id$} (R2);
\end{tikzpicture}
\qquad\qquad
\begin{tikzpicture}[on grid,auto,node distance = 1.5cm,baseline=(a1)]
\node (R1) {$R^b$};
\node (P1) [below=2 of R1] {$P^p_j$};
\node (R2) [left=2 of R1] {$R^c$};
\node (P2) [below=2 of R2] {$P^q_k$};

\path
  (R1) edge [->] node [right] {$a^p_j$} (P1)
  (R2) edge [->] node [left] {$a^q_k$} (P2)
  (R1) edge [->] node {$r$} (R2)
  (P2) edge [->] node [below] {$r_{k,j}}} (P1);

\node (S) [right=2 of R1] {$R^c$};
\path (S) edge [->] node {$v$} (R1);
\end{tikzpicture}
\]
%
In particular, also $v;a^p_j=a^q_k;r_{k,j}=v;r;a^q_k;r_{k,j}$, meaning that $v$ equalises the right hand diagram above. Since, moreover, $v;r=\id$ is epi, the right hand diagram is a (simple) reflection square.

\smallskip
By the induction hypothesis, for all $j\in[1,|b|]$ and $k\in[1,|c|]$, since $m_{j,k}\of c^p_j\to c^q_k$ is a preservation morphism with top arrow $v_{m_{j,k}}$ is a split mono there is a reflection morphism $m'_{k,j}\of c^q_k\to c^p_j$ such that $m_{j,k};m'_{k,j}=\id_{c^p_j}$. It follows that $m'=(r,o^{-1},\setof{m'_{k,j}}_{(k,j)\in o^{-1}})$ is a reflection morphism from $c$ to $b$ (taking into account that $o$ is total and univalent, hence $o^{-1}$ is surjective and injective); and moreover, $m;m'=\id_b$.\qed 
\end{description}
\end{proof}


\todo[inline]{From here is the previous version, based on shifters.}
%
It also implies that morphisms should only relate conditions with the same root, because entailment is defined only among conditions over the same root (see Def.~\ref{def:ab-satisfaction}).\todo{AC: rephrase} Technically, the main result of this section will be formulated categorically as the existence of  contravariant functors from  categories of conditions with structural morphisms to $\ABC^{\entails}$.

The underlying principle for a morphism $m$ from $c=(R,p^c_1\ccdots p^c_{|c|})$  to $b=(R,p^b_1\ccdots p^b_{|b|})$  is that it must identify, for every $b$-branch $p^b_i$, a $c$-branch $p^c_j$ that it entails. Entailment of $p^c_j$ by $p^b_i$ is captured by the existence of an arrow $v_i$ from $P^c_j$ to $P^b_i$  such that there is (recursively, but in the opposite direction with respect to $m$) a morphism between the subconditions $b_i$ and $c_j$; however, this sub-morphism can only exist ``modulo" the arrow $v_i$ between their roots. 
The task of the next subsection is to make the concept of ``modulo $v_i$" precise, introducing the notion of \emph{shifter}.

\subsection{Condition shifting}
\slabel{shift-morphisms}

Given two conditions $c\in \AC{A}$ and $b\in \AC{B}$ and an arrow $v\of A \to B$, we have two main ways to reduce both to the same root, using Def.~\ref{def:shift}. We can either precompose $b$ with $v$ obtaining $\bshift{b}{v}\in \AC{A}$, or we can shift $c$ along $v$, obtaining $\shift{c}{v}\in \AC{B}$. As commented earlier, precomposition is a local operation that only affects the top-level branches of a condition, while shifting is a global operation that affects the branches of a condition at any depth. In the definition of structural morphisms between ab-conditions the roots of subconditions have to be reconciled at each level: this makes the use of shift quite complex to manage. Therefore we will stick to precomposition.\todo{AC: add discussion on shift-related morphisms, when the relationship with span-based ones is more clear, or leave their study to future work.} 

Taking into account the goal of defining morphisms between conditions that witness entailment in the opposite direction, we introduce two operations on conditions, both based on precomposition. 
Given arrow $v:A\to B$, we consider precomposition with $v$ as a function $\cB_v$ from $\AC{B}$ to $\AC{A}$, but only if $v$ is epi, since otherwise we cannot guarantee that $v$-prefixed models of $\bshift{b}{v}$  induce models of $b$, a technical condition needed later on. If $v$ satisfies the stronger condition of being \emph{split epi}, meaning that it has a \emph{section} $s : B \to A$ such that $s;v= id_B$, then we can also precompose with $s$, obtaining a function $\cF^s_v$ from $\AC{A}$ to $\AC{B}$. Let us first recall some categorical terminology.

\begin{definition}[split epi, split mono, retraction, section]\label{def:retraction}
  An arrow $r\of A \to B$  is a \emph{split epi(morphism)} if there is an arrow $s\of B\to A$ such that $s;r = id_B$. In this case, arrow $s$ which satisfies the dual condition is a 
\emph{split mono(morphism)}. We also say that $r$ is a \emph{retraction} of $s$, and that $s$ is a \emph{section} of $r$.
\end{definition} 



% % We call this operation \emph{backward shifting}. If $v$ is split epi, we can also precompose with a section of $v$, obtaining an operation from $\AC{A}$ to $\AC{B}$, which we call \emph{forward shifting}. The existence of a section allows to shift a condition backwards along the section, thus moving the condition forwards along $v$. 


% the backward shifting 

% We introduce two operations on 

% - v must be epi
% - if it is also split-epi, suing a section we can go the other way around
% - being split-epi means that the source obiect can be "folded" to a subobject. For example it can consist of several copies of the suboject, or items can be merged/fused. 
% Having a section allows to shift a condition backwards slong the section, thus moving the condition forwards along v.


% . 
%  Our purpose is to relate the models of a condition $b\in \AC A$ to those of a condition $c\in \AC B$, given an arrow $v:A\to B$ between their roots (the existence of which implies that the model sets are typically disjoint). To be more precise, we will require that \emph{every $v$-prefixed model of $b$ induces a model for $c$}: that is, if $v;g$ is a model of $b$, then $g$ is a model of $c$.



% Let us introduce some categorical terminology.





% Nevertheless, note that if in the above setting $v$ has a (possibly partial) inverse $x\of B \to A$, we may also consider precomposition of $b$ with $x$, obtaining $\bshift{b}{x}\in \AC{A}$. 




% We are interested in the situation where $v$ is a split epi, and hence it has a section $x$. 

% \subsection{Root shifting}
% \slabel{shift-morphisms}

% In order to relate conditions at different roots, we introduce a special kind of mapping called a \emph{root shifter}, which will later turn out to be a functor between categories of conditions over different roots. For now, we note that the 

% \todo{AC: check this****} Our purpose is to relate the models of a condition $b\in \AC A$ to those of a condition $c\in \AC B$, given an arrow $v:A\to B$ between their roots (the existence of which implies that the model sets are typically disjoint). To be more precise, we will require that \emph{every $v$-prefixed model of $b$ induces a model for $c$}: that is, if $v;g$ is a model of $b$, then $g$ is a model of $c$.

% Note: This requires $v$ to be epi, and next that we inductiuvely relate two subconditions havong the same root.  
% This can be achieved in essentially two ways, both on the basis of $v$: by precomposing $b$ with $v$, that  we will call \emph{backward shifting}, or by precomposing $c$ with a section of $v$, that we will call \emph{forward shifting}. It is convenient to introduce a name for these operations, which will be used in the definition of morphism between conditions.

We introduce the two operations sketched above formally as follows.

\begin{definition}[forward and backward shifters]\label{def:ab-shifter}
Given an epi $v\of A\to B$, the \emph{backward shifter for $v$} is the operation $\cB_v: \AC{B} \to \AC{A}$  defined as $\cB_v(b) = \bshift{b}{v}$.
Similarly, if $v$ is split epi, for any section $s$ of $v$ we call a \emph{forward shifter for $v$} the operation $\cF^s_v: \AC{A} \to \AC{B}$ defined as $\cF^s_v(c) = \bshift{c}{s}$.
\end{definition}
\fcite{ex-ab-shifters} shows some examples.\todo{AC: edit the figure to show the shifters acting on conditions, not on arrows.}


\begin{figure}
\centering
\input{figs/ex-ab-shifters}
\caption{Two different shifters $\cF_v^{s_1}$ and $\cF_v^{s_2}$ for $v \of A \to B$ applied to condition $c$, and backward shifter $\cB_{v'}$ for $v'$ applied to $b$. (All morphisms $a$ are the identity on nodes.)}
\flabel{ex-ab-shifters}
\end{figure}

The followings are the crucial preservation properties of shifters:
%
\begin{proposition}[shifters preserve models]\plabel{ab-root shifters preserve}
Let $v:A\to B$.
\begin{enumerate}[topsep=\smallskipamount]
\item If $c\in \AC B$ and $\cB$ is the backward  shifter for $v$, then $v;g\sat \cB(c)$ implies $g\sat c$.
\item If $c\in \AC A$ and $\cF^s_v$ is a forward shifter for $v$, then $v;g\sat c$ implies $g\sat \cF^s_v(c)$.
\end{enumerate}
\end{proposition}
%
\begin{fullorname}
\begin{proof}
Let $v:A\to B$.
\begin{enumerate}
\item Since $\cB_v(c)= \bshift{c}{v}$, the statement follows by clause 2 of \pcite{shiftSat} because $v$ is epi.
\item Recall that $\cF^s_v(c) = \bshift{c}{s}$. By clause 2 of \pcite{shiftSat} we have that $v;g \sat c$ implies $s;v;g \sat \bshift{c}{s}$, hence $g\sat \bshift{c}{s}$ because $s$ is a section of $v$.\qed
\end{enumerate}
\end{proof}
\end{fullorname}
% 
Shifters compose, in the following way.
%
\begin{proposition}[identities and composition of shifters]\plabel{ab-shifters compose}
\begin{enumerate}[topsep=\smallskipamount]
\item For every object $X$, the identity function on $\AC{X}$ is both a backward shifter for $id_X$, $\cB_{id_X}$, and a forward shifter for $id_X$, $\cF^{id_X}_{id_X}$.
% \item If $\cU$ is a source [root] shifter from $X$ to $Y$ and $\cV$ a source [root] shifter from $Y$ to $Z$, then $\cU;\cV$ is a source [root] shifter from $X$ to $Z$. 
\item If $\cB_u$ and $\cB_v$ are backward  shifters for $u: Y\to X$ and $v: Z \to Y$, respectively, then $\cB_u;\cB_v$ is a backward shifter for $v;u: Z \to X$.
\item If $\cF^{s_1}_u$ and $\cF^{s_2}_v$ are forward shifters for $u: X\to Y$ and $v: Y \to Z$, respectively, then $\cF^{s_1}_u;\cF^{s_2}_v$ is a forward shifter for $u;v: X \to Z$.
\item Composition of backward and forward shifters is associative and satisfies the identity laws.
\end{enumerate}
\end{proposition}

\begin{fullorname}
\begin{proof}
\begin{enumerate}
\item Both $\cB_{id_X}$ and $\cF^{id_X}_{id_X}$ are defined ($id_X$ is split epi) and  precompose a condition over $X$ with $id_X$, therefore they act as the identity function on $\AC{X}$. 
\item By associativity of arrow composition and since epi compose, we have  $\cB_u;\cB_v = \cB_{v;u}$.
\item By associativity of arrow composition and since split monos compose, we have $\cF^{s_1}_u;\cF^{s_2}_v = \cF^{s_2;s_1}_{u;v}$.
\item Obvious. \qed
\end{enumerate}
\end{proof}
\end{fullorname}

In the definition of structural morphisms in the next section backward and forward shifters will be used to relate (sub)conditions over different roots, exploiting the preservation of models (\pcite{ab-root shifters preserve}). 

% However, shifters require the arrow relating the roots to be (split) epi, which could be too restrictive.  
% There is a border situation where two conditions over different roots related by an \emph{arbitrary} arrow satisfy the model preservation property: this is the case when the first condition has depth 0, and the property holds vacuously.

% \begin{fact}[zero shifters]
%   \flabel{zero shifters}
% Let $v: A \to B$ be an arrow and $b \in \AC{B}$ be an ab-condition. Then $v;g\sat (A,\epsilon)$ implies $g \sat b$.
% %
% This is obvious because conditions of depth 0 have no models. 

% % We will need a name for this situation, thus we call a \emph{zero shifter} the pair $(v,b)$, intended to relate $b$ with $(A,\epsilon)$.\todo{AC: possibly delete if not used anymore.}
% \end{fact}

\subsection{Structural morphisms}

The following defines structural morphisms, based on shifters, over ab-conditions with the same root.
 
\begin{definition}[arrow-based condition morphism]\label{def:ab-morphism}
  Let $c,b \in \AC{R}$. A \emph{forward-shift [backward-shift] ab-condition morphism (or ab-morphism)} $m: c \to b$ is a pair $(o,(v_1,m_1)\ccdots (v_{|b|},m_{|b|}))$ where
  \begin{enumerate}[topsep=\smallskipamount]
  \item $o:[1,|b|]\to[1,|c|]$ is a function from $b$'s branches to $c$'s branches; 
  \item for all $1\leq i\leq |b|$, $v_i:P^c_{o(i)}\to P^b_i$ is an arrow from the pattern of $p^c_{o(i)}$ to that of $p^b_i$ such that $a^c_{o(i)};v_i=a^b_i$;
  \item\emph{Forward shift}: for all $1\leq i\leq |b|$, either $c_{o(i)}$ has depth 0 and 
  $m_i: b_i \to \zeroCond{P^b_i}$ is the  
  terminal
  morphism, 
  %$m_i = (v_i,b_i)$ is a zero shifter, 
  or there is a forward  shifter $\cF_i$ for $v_i$ such that $m_i:b_i \to \cF_i(c_{o(i)})$ is a forward-shift morphism.
  \item\emph{Backward shift}: for all $1\leq i\leq |b|$, either $c_{o(i)}$ has depth 0 and 
  %$m_i = (v_i,b_i)$ is a zero shifter, 
  $m_i: \bshift{b_i}{v_i} \to c_{o(i)} (= \zeroCond{P^c_{o(i)}})$ is the 
  terminal 
  morphism,
  or there is a backward  shifter $\cB_i$ for $v_i$ such that $m_i:\cB_i(b_i) \to c_{o(i)}$ is a backward-shift morphism.
  \end{enumerate}
\end{definition}
 %
We use the same notational conventions as for conditions: $m$ has width $|m|$ and depth $\depth(m)$, and the components of a morphism $m$ are denoted $o^m$, $v^m_i$ and $m_i$ for all $1\leq i\leq |m|$. 

The special cases involving conditions of depth 0 in clauses 3 and 4 of the previous definition are not strictly necessary, but exploting the fact that such conditions have no models (Fact~\ref{fact:satisfaction}) they define a richer class of morphisms.
Morphisms involving conditions of depth 0 enjoy some relevant properties, that we will exploit later. We also show that the disjunctions in clauses 3 and 4 of Def.~\ref{def:ab-morphism} do not introduce ambiguity. 

\begin{lemma}[terminal morphisms and 0-depth conditions]
             \llabel{terminal morphisms}
\begin{enumerate}
\item
For each ab-condition $b \in \AC{A}$ there is a unique \emph{terminal morphism} $!_b: b \to (A,\epsilon)$ to the 0-depth condition over $A$. 
\item 
Given a morphism of ab-conditions $m:c \to b$, if $c$ has depth 0 then also $b$ has depth 0, thus $m =\, !_c = id_c$. 
\item 
The disjunction in clause 3  of Def.~\ref{def:ab-morphism} is not ambiguous, because if $c_{o(i)}$ has depth 0 and there is a forward shifter $\cF_i$ for $v_i$,  then 
$m_i: b_i \to \cF_i(c_{o(i)})$ is the terminal morphism.
\item 
The disjunction in clause 4  of Def.~\ref{def:ab-morphism} is not ambiguous, because if $c_{o(i)}$ has depth 0 and there is a backward shifter $\cB_i$ for $v_i$, then $m_i$ is the terminal morphism $!_{\bshift{b_i}{v_i}}: \bshift{b_i}{v_i} \to c_{o(i)}$.
\end{enumerate}
\end{lemma}


\begin{proof}
\begin{enumerate}
  \item For clause 1 of Def.~\ref{def:ab-morphism} there is a unique function $o: \varnothing \to [1,|b|]$, and the other clauses of are vacuous.
  \item If $c$ has depth 0 set $[1, |c|]$ is empty, thus the function $o$ of clause 1 of Def.~\ref{def:ab-morphism} is well defined only if $|b| = 0$, which also implies that $b=c$.
  \item Observe that forward shifters preserve depth, because they are defined in terms of precomposition  with a section. Therefore $\cF_i(c_{o(i)})$ has depth 0.
\item  Simply recall that if $\cB_i$ is the backward shifter for $v_i$ then $\cB_i(b_i) = \bshift{b_i}{v_i}$ \qed 
\end{enumerate}
\end{proof}
%
\fcite{shifted-morphism} visualises the principle of forward- and backward-shift morphisms.
%
\begin{figure}[t]
\centering
\input{figs/shifted-morphisms}
\caption{Visualisation of forward-shift and backward-shift morphisms $m:c \to b$}
\flabel{shifted-morphism}
\end{figure}

\begin{example}\exlabel{ab-morphisms}
\fcite{ab-conditions-morphisms} depicts a morphism from the third to the second ab-conditions of \fcite{ab-conditions}. We use here plain labels (like $c_2, c_{321}$) to denote graphs in the picture, and hatted labels (like $\widehat{c_2}, \widehat{c_{321}}$) for the conditions rooted at them. 
Then  there is a backward-shift morphism $m\of \widehat{c_3}\to \widehat{c_2}$ defined as $m = (\{1\mapsto 2\}, (v_1: c_{32} \to c_{21}, m'))$, where $m':\cB_{v_1}(\widehat{c_{21}}) \to \widehat{c_{32}}$ is 
$m' = (\{1\mapsto 2\},(v_{12}: c_{212} \to c_{321}, m''))$, and $m'':\bshift{\widehat{c_{321}}}{v_{12}} \to \widehat{c_{212}}$ is the terminal morphism to the 0-depth condition $\widehat{c_{212}}$. 
% $m' = (\{1\mapsto 2\},(v_{12}: c_{212} \to c_{321}, (v_{12}, \widehat{c_{321}})))$. 

Actually, since $v_1$ is the identity there is also a forward-shift morphism 
$m_1\of \widehat{c_3}\to \widehat{c_2}$, with $m_1 = (\{1\mapsto 2\}, (v_1: c_{32} \to c_{21}, m'))$, and 
$m':\widehat{c_{21}} \to \cF_{v_1}^{v_1}(\widehat{c_{32}})$ defined as 
$m' = (\{1\mapsto 2\},(v_{12}: c_{212} \to c_{321}, m''))$, and $m'':\widehat{c_{321}} \to \zeroCond{c_{321}}$ is the terminal morphism. 
\qed
\end{example}

\begin{figure}[t]
\centering
\input{figs/ex-ab-conditions-morphism}
\caption{Example of morphism between arrow-based conditions (see \excite{ab-morphisms})}
\flabel{ab-conditions-morphisms}
\end{figure}



As we shall see in \pcite{ab-morphisms preserve models} these notions of morphism witness entailment in the opposite directions, but let us first show that they turn $\AC{R}$ into a category. To this aim,  we need to show that morphisms compose. This, in turn, depends on the preservation of morphisms by shifters.
 
\begin{lemma}[shifters preserve morphisms]
             \llabel{ab-root shifters preserve morphisms}
Let $c,b \in  \AC{X}$ and let $v\of Y\to X$. If $m:c\to b$ is a morphism, then also $m:\bshift{c}{v}\to \bshift{b}{v}$ is a morphism.

As a consequence if $\cB_v$ is the backward shifter for $v\of Y\to X$. If $m:c\to b$, then also $m:\cB_v(c)\to\cB_v(b)$. 
Similarly, if $\cF^s_v$ is a forward shifter for $v\of X\to Z$, if $m\of c\to b$, then also $m:\cF^s_v(c)\to\cF^s_v(b)$. 

% Let $c,b \in  \AC{X}$ and let $\cC$ be a root shifter from $X$ to $Y$. If $m:c\to b$, then also $m:\cC(c)\to\cC(b)$. 
\end{lemma}
%
\begin{fullorname}
\begin{proof}
Assume $m = (o,(v_1,m_1)\cdots(v_{|b|},m_{|b|}))$. The only new proof obligation for $m$ to be an ab-morphism from $\bshift{c}{v}\to \bshift{b}{v}$ is that, for all $1\leq i\leq |b|$, $v;a^c_{o(i)};v_i=v;a^b_i$. This follows immediately from $a^c_{o(i)};v_i=a^b_i$. 

The other statements follow immediately from the first one, since both backward and forward shifters are defined in terms of backward shift along an arrow (see Def.~\ref{def:ab-shifter}).\todo{AC: this means that we can change the source of a morphism $m$, if needed. This is what we do for composing morphism, applying a shifter to match the target of a morphism with the source of another. We should do the same for the identity of the empty condition, to shift it along an arrow. This was the trivial shifter. Can we call these morphisms polymorphic? After all the identity on the zero depth condition is just $id_{\zeroCond{A}} = (o = \varnothing, \epsilon)$}
%$\cB_v(b) = \bshift{b}{v}$ and $\cF^s_v(c) = \bshift{c}{s}$.
%A similar observation holds for forward shifters, considering the section $s\of Z \to X$ instead of $v$. \qed
\end{proof}
\end{fullorname}
%
%


Composition and identities of ab-condition morphisms can be defined inductively, and are independent of their forward or backward nature.  

\begin{definition}[identities and composition of ab-morphisms]\label{def:ab-morphism-composition}
Let $e,c,b$ be ab-conditions and $m:e\to c,n:c\to b$ ab-morphisms, with $|b|=w$. Then $\id_b$ and $m;n$ are defined inductively as follows:
%
\begin{eqnarray}
\id_b & =
  & (\id_{[1,w]},(\id_{P^b_1},\id_{b_1})\cdots 
                 (\id_{P^b_w},\id_{b_w}))
  \label{eq:id-morphism} \\
m;n & =
  & (o^n;o^m,(v^m_{o^n(1)};v^n_1,n_1;m_{o^n(1)})\cdots 
              (v^m_{o^n(w)};v^n_w,n_w;m_{o^n(w)})) \enspace.
%               & (o^m;o^n,(v^n_{o^m(1)};v^m_1,n_{o^m(1)};m_1)\cdots 
%               (v^n_{o^m(w)};v^m_w,n_{o^m(w)};m_w)) \enspace.
 \label{eq:morphism composition}
\end{eqnarray}

\end{definition}
%
Identities are morphisms, and morphism composition is well-defined.

\begin{lemma}[identities and composition are well defined]\llabel{ab-morphisms compose}
Let $e,c,b$ be ab-conditions and $m\of e\to c, n\of c\to b$ ab-condition morphisms.
\begin{enumerate}[topsep=\smallskipamount]
\item $\id_b$ is both a forward- and a backward-shift morphism from $b$ to $b$;
\item If $m$ and $n$ are forward-shift morphisms, then $m;n$ is a forward-shift morphism from $e$ to $b$;
\item If $m$ and $n$ are backward-shift morphisms, then $m;n$ is a backward-shift morphism from $e$ to $b$.
\end{enumerate}
\end{lemma}
%
\begin{fullorname}
\begin{proof}
Let $w=|b|$.
%
\begin{enumerate}[topsep=\smallskipamount]
\item By induction on the depth of $b$. For all $1\leq i\leq w$, $a^b_i;\id_{P^b_i}=a^b_i$. Recall by clause 1 of \pcite{ab-shifters compose} that $\rF{\id}{\id}$ and $\cB_\id$ (with $\id=\id_{P^b_i}$) are, respectively, forward and backward shifters, both acting as the identity on arrows with source $P^b_i$; hence (by the induction hypothesis) $\id_{b_i}$ is both a forward-shift morphism from $b_i$ to $\rF\id\id(b_i)$ and a backward-shift morphism from  $\cB_\id(b_i)$ to $b_i$.
\end{enumerate}
%
The other two clauses are proved by induction on the depth of $m$ and $n$. For clause 1 of Def.~\ref{def:ab-morphism}, let $o=o^n;o^m$, which is a function from $[1,w]$ to $[1,|e|]$ as needed. For clause 2, let $v_i=v^m_{o^n(i)};v^n_i$ for all $1\leq i\leq w$. All $v_i$ are arrows from $P^e_{o(i)}$ to $P^b_i$, and they commute with the arrows from the root, as required. In fact, $a^e_{o(i)};v_i =\text{(def.~of $v_i$) } a^e_{o(i)};v^m_{o^n(i)};v^n_i  =\text{($m$ is a morphism) }  a^c_{o^n(i)} ; v^n_i =\text{($n$ is a morphism) }  a^b_i$.
%
\begin{enumerate}[resume]
\item
For clause 3 of Def.~\ref{def:ab-morphism} we have to show that for all $1\leq i\leq w$, either $e_{o(i)}$ has depth 0 and $n_i;m_{o^n(i)}: b_i \to \zeroCond{P^b_i}$ is the terminal morphism, or there is a forward root shifter $\cF_i$ for $v_i$ such that $n_i;m_{o^n(i)}:b_i \to \cF_i(e_{o(i)})$ is a forward-shift morphism. 

Let $j = o^n(i)$ and let us proceed by cases. If both $n_i:b_i \to \cF^n_i(c_{j})$ and  $m_{j}:c_j \to \cF^m_j(e_{o^m(j)})$ are forward shift morphisms, we have $m_{j}:\cF^n_i(c_j) \to \cF^n_i(\cF^m_j(e_{o^m(j)}))$ by \lcite{ab-root shifters preserve morphisms}, and hence $\cF^m_j;\cF^n_i$ is a forward shifter for $v_i$ by \pcite{ab-shifters compose} and, using the inductive hypothesis, $n_i;m_{o^n(i)}:b_i \to \cF^n_i(\cF^m_j(e_{o(i)}))$  is a forward shift morphism. 

If $e_{o(i)}$ has depth 0, then $m_{j}:c_j \to \zeroCond{P^c_j}$ is the terminal morphims, by Def.~\ref{def:ab-morphism} and clause 3 of \lcite{terminal morphisms}. If $n_i:b_i \to \cF^n_i(c_{j})$ is a forward shift morphism then we obtain 
$n_i;m_{o^n(i)}:b_i \to \cF^n_i(\zeroCond{P^c_j})$ as above, and observe that $\cF^n_i(\zeroCond{P^c_j}) = \zeroCond{P^b_i}$, as desired. Otherwise, if $c_j$ has depth 0, then
$n_i: b_i \to \zeroCond{P^b_i}$ is the terminal morphism, and  $m_{j}:c_j \to \zeroCond{P^c_j}$ is the identity defined as $m_{j} = id_{\zeroCond{P^c_j}} = (o = \varnothing, \epsilon)$, that we can see as well as $id_{\zeroCond{P^b_i}}$; thus $n_i;m_{o^n(i)}: b_i \to \zeroCond{P^b_i}$ is the terminal morphism, as required. 

Finally observe that if $c_j$ has depth 0 and thus
$n_i: b_i \to \zeroCond{P^b_i}$ is the terminal morphism,
then $e_{o(i)}$ has depth 0 as well, and we fall in the previous cases. In fact, otherwise $m_{o^n(i)}: c_j \to \cF^m_j(e_{o(i)})$ should be a forward shift morphism, but by clause 2 of \lcite{terminal morphisms} $\cF^m_j(e_{o(i)})$ should have depth 0, and since shifters preserve depths  $e_{o(i)}$ should have depth 0 as well. 


\item Mutatis mutandis.
\qed
\end{enumerate}
\end{proof}
\end{fullorname}
%

The associativity of morphism composition and the identity laws can be proved by induction on the depth of the morphisms, using the associativity and identity laws for shifters and morphisms at lower depth. We omit the details. This ensures that for each object $R$ the conditions over $R$ with either forward-shift or backward-shift morphisms give rise to a well-defined category. Moreover, since there are no morphisms between conditions with different roots, we obtain two categories $\ABC^\forw$ and $\ABC^\back$ with ab-conditions as objects and forward- and backward-shift morphisms as arrows, respectively.\todo{AC: More natural to define the full categories first, and then the fibres as full subcategories. }

\iffull
\begin{definition}[categories of ab-conditions]
  \label{def:ab-morphism categories}
For every $R$, $\ABC^\forw(R)$  is the category having all conditions over $R$ as objects and  forward-shift ab-morphisms as arrows.  Similarly, $\ABC^\back(R)$ has $\AC R$ as objects and backward-shift ab-morphisms as arrows.

Furthermore, $\ABC^{\forw}$ is the category having all ab-conditions as objects and forward-shift morphisms as arrows, and  $\ABC^{\back}$ is the category having the same objects and backward-shift morphisms as arrows. \todo{AC: Add explicit proof of well-definedness?}
\end{definition}
%
\fi


The next result is the main motivation for the introduction of structural morphism among conditions. It states that the existence of a morphism between two conditions $m: c \to b$ provides evidence that $b\entails c$.

\begin{proposition}[ab-condition morphisms reflect models]
  \plabel{ab-morphisms preserve models}
Let $c,b \in \AC{R}$  be arrow-based conditions. If $m:c\to b$ is a backward-shift morphism or a forward-shift one, then $g\sat b$ implies $g \sat c$ for all arrows $g:R\to G$. 
\end{proposition}
%
\begin{fullorname}
\begin{proof}
Let $m = (o,(v_1,m_1)\ccdots (v_{|b|},m_{|b|}))$ and assume by induction hypothesis that for all $1\leq j\leq |b|$ the property holds for $m_j$. Let $p^b_i$ be the responsible branch and $h: P^b_i \to G$ the witness for $g\sat b$. 
Hence we have $(\dagger)\, g=a^b_i;h$ and $(\ddagger)\, h \nsat b_i$.  
We show that $g \sat c$, with responsible branch $p^c_{o(i)}$ and witness $v_i; h$. In fact, we immediately have  $a^c_{o(i)}; v_i ;h = a^b_i ; h \stackrel{(\dagger)}{=} g$. It remains to show that $v_i ;h \nsat c_{o(i)}$, for which we consider separately the two cases.
\begin{enumerate}
% \item Let $m$ be a forward-shift morphism. Then for all $1\leq j\leq |b|$ we have that   $m_j:b_j \to \cF_j(c_{o(j)})$ is a forward-shift morphism, where $\cF_j$ is a forward shifter for $v_j$. If (ad absurdum) $v_i ;h \sat c_{o(i)}$, by clause 2 of \pcite{ab-root shifters preserve} we would have $h \sat \cF_i(c_{o(i)})$, and by induction hypothesis on $m_i$ also $h \sat b_i$, contradicting $(\ddagger)$.
\item Let $m$ be a forward-shift morphism. Then for all $1\leq j\leq |b|$ we have that  either $c_{o(j)}$ has depth 0 (and  $m_j:b_j \to \zeroCond{P^b_j}$ is the terminal morphism), or  $m_j:b_j \to \cF_j(c_{o(j)})$ is a forward-shift morphism, where $\cF_j$ is a forward shifter for $v_j$. 

If $c_{o(i)}$ has depth 0 then it has no models (Fact~\ref{fact:satisfaction}), thus $v_i ;h \not \sat c_{o(i)}$. Otherwise, if (ad absurdum) $v_i ;h \sat c_{o(i)}$, by clause 2 of \pcite{ab-root shifters preserve} we would have $h \sat \cF_i(c_{o(i)})$, and by induction hypothesis on $m_i$ also $h \sat b_i$, contradicting $(\ddagger)$.

\item Let $m$ be a backward-shift morphism. Then for all $1\leq j\leq |b|$ we have that  either $c_{o(j)}$ has depth 0 (and $m_j:\bshift{b_j}{v_j}\to c_{o(j)}$ is the terminal morphism), or $m_j:\cB_j(b_j) \to c_{o(j)}$ is a backward-shift morphism, where $\cB_j$ is the backward  shifter for $v_j$.

If $c_{o(i)}$ has depth 0 then it has no models, thus $v_i ;h \not \sat c_{o(i)}$.
 Otherwise, suppose (ad absurdum) that $v_i ;h \sat c_{o(i)}$. Then by the induction hypothesis on $m_i$ we also have $v_i ;h \sat \cB_i(b_i)$, and by clause 1 of \pcite{ab-root shifters preserve} we  have $h \sat b_i$, contradicting $(\ddagger)$.

% \item Let $m$ be a backward-shift morphism. Then for all $1\leq j\leq |b|$ we have that  $m_j:\cB_j(b_j) \to c_{o(j)}$ is a backward-shift morphism, where $\cB_j$ is a backward  shifter of width $|b_j|$.
% Suppose (ad absurdum) that $v_i ;h \sat c_{o(i)}$. Then by the induction hypothesis on $m_i$ we also have $v_i ;h \sat \cB_i(b_i)$, and by clause 1 of \pcite{ab-root shifters preserve} we  have $h \sat b_i$, contradicting $(\ddagger)$.

\qed
\end{enumerate}
\end{proof}
\end{fullorname}

\begin{example}\exlabel{ab-morphisms-1}
Going back to \fcite{ab-conditions}, though (as noted before) $c_1\entails c_2$, there does not exist a morphism $m\of c_2\to c_1$. A mapping $v_1$ does exist from $c_{21}$ to $c_{11}$; however, no mapping $v_{11}$ can be found to $c_{211}$ from either $c_{111}$ or $c_{112}$. 
%, nor a mapping $v_{12}$ to $c_{212}$.

On the other hand as shown in \excite{ab-morphisms}
$c_2\entails c_3$ has evidence in the form of both a forward-shift and a backward-shift morphism from $c_3\to c_2$. \todo{Add example showing that there are obvious morphisms that would not exist if 0 depth conditions are not handled explicitely: A => A and B, A => (A and B) and not C ***}
\end{example}
%
Proposition~\ref{prop:ab-morphisms preserve models} can be phrased more abstractly at the categorical level: our first main result is that ab-conditions endowed with either forward or backward morphisms give rise to a category with a contravariant functor to $\ABC^{\entails}$.

\begin{theorem}[categories of arrow-based conditions]\thlabel{ab-categories}
%Let $\ABC^{\forw}$ be the category having ab-conditions as objects and forward-shift morphisms as arrows, and let $\ABC^{\back}$ be the category having the same objects and backward-shift morphisms as arrows. Both categories are well-defined, and there are identity-on-objects functors $\op{(\ABC^{\forw})} \to \ABC^{\entails}$ and $\op{(\ABC^{\back})} \to \ABC^{\entails}$.
Let $\ABC^{\forw}$  and  $\ABC^{\back}$ be the categories defined in Def.~\ref{def:ab-morphism categories}. There are identity-on-objects functors $\op{(\ABC^{\forw})} \to \ABC^{\entails}$ and $\op{(\ABC^{\back})} \to \ABC^{\entails}$.
\end{theorem} 
%
\emph{Proof sketch.} \todo{AC:this is more about well-definedness of the categories.} Composition of morphisms and the fact that all identity morphisms have the correct nature for both categories is shown by \lcite{ab-morphisms compose}.
Moreover, identity morphisms and morphism composition satisfy the unit and associativity laws. The existence of the two identity-on-object functors follows from \pcite{ab-morphisms preserve models}.
\qed 

\medskip\noindent
Even though both forward- and backward-shift morphisms reflect models, they are incomparable: if two conditions are related by one kind of morphisms they are not necessarily related by the other kind. In other words, if we consider the existence of a morphism as providing evidence for entailment, then these two types of morphism are complementary, as the following example shows.
%
\begin{figure}[t]
\centering
\subcaptionbox
  {A forward-shift morphism $m \of c \to b$ without a backward-shift one
   \flabel{forward-no-backward}
  }
  [.45\textwidth]
  {\input{figs/ex-forward-no-backward}}
  \qquad
\subcaptionbox
  {A backward-shift morphism $m \of c \to b$ without a forward-shift one
   \flabel{backward-no-forward}
  }
  [.45\textwidth]
  {\input{figs/ex-backward-no-forward}}
\caption{Forward-shift versus backward-shift morphisms}
\flabel{forward vs backward}
\end{figure}
%
\begin{example}[forward versus backward root shifters]\exlabel{forward vs backward}
\fcite{forward-no-backward} shows a forward-shift morphism $m \of c \to b$. 
% where there is no backward-shift one. 
Both conditions are evaluated in a context where we know $\la(x,z)$. In that context,  $c$ is equivalent to $\exists y\st \la(x,y) \wedge \neg \lb(z,x)$, and $b$ is equivalent to the property $\neg \lb(z,x)$. Intuitively, $b\entails c$ holds because $y$ can be instantiated with $z$. The figure shows $\cF_{m_1}^s(a^c_{11})$ (for a section $s$ of $m_1$), which acts as the identity on node and edges. Note that there is no backward-shift morphism $m'\of c \to b$ because $\cB_{m_1}(a^b_{11}) ; m_{11} = m_1; a^b_{11} ; m_{11} \not = a^c_{11}$.

\fcite{backward-no-forward} shows a backward-shift morphism $m \of c \to b$. 
% where there is no forward-shift one.
Both conditions are evaluated in a context where we know $\la(x,y)$. In that context, $c$ is equivalent to $\exists z.\lb(z,z) \wedge \neg(y=z \wedge \lc(y,x))$ and $b$ is equivalent to $\lb(y,y) \wedge \neg \exists z\st \lc(z,x)$.
Indeed, $b\entails c$ because $z$ can be instantiated with the value of $y$. The figure shows $\cB_{m_1}(a^b_{11})$. Note that there cannot be a forward-shift morphism $m'\of c \to b$ because $m_1$ is not split epi.
\qed
\end{example}
%
As a final observation,\todo{AC:to be checked.} we can characterise the relation between shifters and conditions as \emph{indexed categories}. For this purpose, let $\bC^\ashift$ denote the category with objects from $\bC$ (which as usual can be any presheaf topos) and arrows $\cS\of X\to Y$ whenever $\cS$ is an (arrow-based) source shifter from $Y$ to $X$. Identity and composition are defined as in \pcite{ab-shifters compose}.

Moreover, for every $\bC^\ashift$-object $X$ let $\ACx^\forw(X)=\ABC^\forw(X)$ (the full subcategory of $\ABC^\forw$ containing only conditions rooted at $X$), and for every $\bC^\ashift$-arrow $\cS\of X\to Y$ let $\ACx^\forw(\cS)$ be the functor that maps every $c\in \AC Y$ to $\bar\cS(c)\in \AC X$ and every forward-shift morphism $m\of c\to b$ (for $c,b\in \AC Y$) to $m\of \bar\cS(c)\to \bar\cS(b)$. Define $\ACx^\back$ analogously for backward-shift morphisms.

\begin{proposition}[$\bC^\ashift$-indexed categories]\plabel{ab indexed}
$\ACx^\forw$ and $\ACx^\back$ are functors from $\op{(\bC^\ashift)}$ to $\Cat$, establishing $\bC^\ashift$-indexed categories.
\end{proposition}
%
Note that this is only correct because (implicitly) $\bC^\ashift$ is a well-defined category, and for all $\cS\of X\to Y$, $\ACx^\forw(\cS)$ and $\ACx^\back(\cS)$ are well-defined functors from $\ABC^\forw(Y)$ to $\ABC^\forw(X)$ and from $\ABC^\back(Y)$ to $\ABC^\back(X)$, respectively.
